\section{曲线：弧长参数化与曲率}
\label{sec:12-Real-Analysis-Support/11-Curve-and-Surface-Geometry/01}

高速公路的匝道从来不是直线接圆弧。真那样修，车开到接缝处，方向盘得在一瞬间从零打到某个固定角度，人做不到，做到了乘客也受不了。设计规范里用的是回旋线，它的曲率随里程线性增长，方向盘可以匀速地转。这里出现了一个需要被量化的东西：``弯得有多急''。它显然不该取决于车开多快——同一段匝道，时速 \(40\) 公里和时速 \(80\) 公里过弯，路的形状一个样。

难点也就在这儿。描述曲线要用参数化，而参数化里混进了``走法''，导数于是同时携带几何信息和无关的走法信息。这一节把走法固定下来，让剩下的导数只说几何。

\begin{center}
\begin{tikzpicture}[>=stealth,scale=0.8]
% 左：缓弯，密切圆大
\draw[dashed,gray!45] (0,2.4) circle (2.4);
\draw[gray!55] (0,2.4) -- (0,0);
\fill[gray!60] (0,2.4) circle (1.6pt);
\node[gray!55!black,font=\small,left=1pt] at (0,1.2) {\(1/\kappa\)};
\draw[orange!85!black,line width=1.6pt] (0,2.4) ++(-128:2.4) arc (-128:-52:2.4);
\fill (0,0) circle (2pt);
\node[font=\small,align=center] at (0,-0.85) {缓弯：贴合的圆大\\曲率 \(\kappa\) 小};
% 右：急弯，密切圆小
\draw[dashed,gray!45] (6.3,0.9) circle (0.9);
\draw[gray!55] (6.3,0.9) -- (6.3,0);
\fill[gray!60] (6.3,0.9) circle (1.6pt);
\node[gray!55!black,font=\small,left=1pt] at (6.3,0.45) {\(1/\kappa\)};
\draw[orange!85!black,line width=1.6pt] (6.3,0.9) ++(-145:0.9) arc (-145:-35:0.9);
\fill (6.3,0) circle (2pt);
\node[font=\small,align=center] at (6.3,-0.85) {急弯：贴合的圆小\\曲率 \(\kappa\) 大};
\end{tikzpicture}
\end{center}

\emph{在每一点贴一个圆，它的半径倒数就是曲率；圆越小，曲率越大。}

\subsection{导数混着几何和走法，得先把走法固定住}

半径 \(R\) 的圆可以这样参数化，也可以那样：

\[
\gamma_{1}(t)=(R\cos t,\,R\sin t),\qquad
\gamma_{2}(t)=(R\cos 2t,\,R\sin 2t).
\]

两者画出同一个圆，但 \(\norm{\gamma_{1}'}=R\)，\(\norm{\gamma_{2}'}=2R\)，加速度的模也差四倍。若拿速度或加速度当``弯曲程度''，同一个圆会得到两个答案。毛病不在圆身上，在参数化身上。

对策直截了当：只准以单位速度走。走法一旦被钉死，导数里剩下的就只有几何。

\subsection{弧长参数化就是以单位速度走完全程}

设 \(\gamma:[a,b]\to\R^{n}\) 光滑且正则，即 \(\gamma'(t)\neq0\) 处处成立。定义弧长

\[
s(t)=\int_{a}^{t}\norm{\gamma'(\tau)}\,d\tau.
\]

正则性使 \(s'(t)=\norm{\gamma'(t)}>0\)，于是 \(s\) 严格递增、可反解出 \(t=t(s)\)。令 \(\widetilde\gamma(s)=\gamma(t(s))\)，链式法则给出

\[
\widetilde\gamma'(s)=\gamma'(t(s))\,t'(s)=\frac{\gamma'(t(s))}{\norm{\gamma'(t(s))}},
\]

模长恒为一。这就是弧长参数化：参数 \(s\) 直接是走过的路程。

正则性这个条件不能省。\(\gamma'(t_{0})=0\) 的点上曲线可能有尖角——想想 \(\gamma(t)=(t^{2},t^{3})\) 在原点，轨迹在那里拐了个尖，而参数化本身光滑得很。正则性挡掉的正是这种``参数光滑但形状不光滑''的情形。

\subsection{曲率是单位切向量转动的速率}

在弧长参数化下，\(T(s)=\widetilde\gamma'(s)\) 是单位切向量，它只能转、不能伸缩。转得多快，就是弯得多急：

\[
\kappa(s)=\norm{T'(s)}=\norm{\widetilde\gamma''(s)}.
\]

由 \(\ip{T}{T}=1\) 求导得 \(\ip{T'}{T}=0\)，所以 \(T'\) 总是垂直于 \(T\)——切向量的变化全部用在转向上，没有一点浪费在改变速率上，这正是把速度固定为一换来的清爽。

两个例子够用了。直线的 \(T\) 恒定，\(\kappa\equiv0\)。半径 \(R\) 的圆用弧长参数化写成 \((R\cos(s/R),R\sin(s/R))\)，二阶导的模是 \(1/R\)，于是 \(\kappa\equiv1/R\)——半径越小弯得越急，这与直觉一致，也说明 \(1/\kappa\) 该被叫做曲率半径。开头那张图画的就是这件事：在曲线的每一点贴上一个与它二阶吻合的圆（密切圆），半径正是 \(1/\kappa\)。匝道设计里的回旋线，说白了就是让 \(\kappa\) 随 \(s\) 线性增加，好让曲率半径从无穷大平滑地收缩到目标值。

实际计算时未必要真去做弧长参数化，一般参数下有现成公式，平面上是

\[
\kappa=\frac{\abs{x'y''-y'x''}}{\bigl((x')^{2}+(y')^{2}\bigr)^{3/2}},
\]

分母上那个 \(3/2\) 次幂就是在替你完成重新参数化。

\subsection{换个走法，曲率纹丝不动}

现在可以兑现开头那句``路的形状一个样''。

\begin{proposition}
设 \(\gamma\) 正则，\(h\) 是正则的参数变换（\(h'\neq0\)），则 \(\gamma\circ h\) 与 \(\gamma\) 在对应点处曲率相同。
\end{proposition}

证明骨架：两条曲线的像集相同，弧长函数只差一个平移或反号，因此它们的弧长参数化至多相差 \(s\mapsto\pm s+c\)。而 \(\kappa=\norm{\widetilde\gamma''}\) 对 \(s\) 的平移显然不敏感，对 \(s\mapsto-s\) 也不敏感——二阶导数遇到符号翻转，两次负号抵消。条件对账：全程只用到 \(h\) 正则（保证 \(\gamma\circ h\) 仍正则、弧长函数仍可反解）与曲线正则本身。

这是本书里第一个明确的``几何量与坐标无关''的例子，值得记住它的模样。曲率不属于任何一种参数化，它属于曲线。上一单元立下的那条纪律——只承认在换坐标后表现良好的量——在这里第一次有了具体的检验对象。后面的第一基本形式、Gauss 曲率，都会重复这套动作：先写出依赖坐标的表达式，再验证它其实不依赖。

\subsection{空间里的曲线除了弯，还能拧}

平面曲线只有一种弯法，向左或向右。到了 \(\R^{3}\)，曲线还能从原来的平面里旋出去。捕捉这一点需要一整套沿曲线活动的坐标系。

设 \(\gamma(s)\) 是空间中的弧长参数化曲线且 \(\kappa>0\)，令 \(T=\gamma'\)，\(N=T'/\kappa\)，\(B=T\times N\)。\(T\) 是切向，\(N\) 指向弯曲的内侧，\(B\) 垂直于 \(T,N\) 张成的密切平面。三者两两正交、长度为一，构成 Frenet 标架。标架沿曲线的变化率写出来是

\[
T'=\kappa N,\qquad
N'=-\kappa T+\tau B,\qquad
B'=-\tau N,
\]

新出现的 \(\tau\) 叫挠率，量的是曲线脱离密切平面的速度。平面曲线的 \(\tau\equiv0\)，因为它从头到尾躺在同一张平面里。圆柱螺线

\[
\gamma(s)=\left(R\cos\frac{s}{c},\,R\sin\frac{s}{c},\,\frac{hs}{c}\right),
\qquad c=\sqrt{R^{2}+h^{2}}
\]

则有常曲率 \(\kappa=R/c^{2}\) 与常挠率 \(\tau=h/c^{2}\)，既弯又拧。

Frenet 方程真正的分量在于它的逆命题：给定连续函数 \(\kappa>0\) 与 \(\tau\)，存在一条曲线以它们为曲率和挠率，并且这条曲线在刚体运动的意义下唯一。两个局部的数完全决定了整体形状。这个模式后面还会遇到——Gauss 曲率对曲面，Fisher 度量对统计流形，走的都是同一条路。

\subsection{一维只能弯一次，二维不行}

曲线在每一点只有一个可弯的方向，一个数就说完了。曲面不一样：站在曲面上的某点，朝不同方向切下去，得到的曲线弯法各不相同——沿一个方向可能鼓起，沿另一个方向可能凹陷。要描述这种局面，一个数远远不够。下一节\hyperref[sec:12-Real-Analysis-Support/11-Curve-and-Surface-Geometry/02]{``曲面：第一与第二基本形式''}引入两个二次型，一个量曲面内部能测到的东西，一个量它在外部空间里怎样弯——而这两者的分野，是本单元最要紧的一件事。
